\label{chap:apendice_exploratorios}

Before the single protocol, three KAN architectures trained from scratch were explored to see how far a KAN could go without a pre-trained feature extractor. They are recorded here for completeness. Their numbers come from a different setting and are \emph{not} comparable with any result of Chapter~\ref{chap:results}.

\begin{itemize}
    \item \textbf{ResKANet-18:} a residual network whose convolutions are themselves parameterised by splines (\texttt{reskanet\_18x32p} from \texttt{torch-conv-kan}~\cite{drokin2024torchconvkan}), trained from scratch with $128 \times 128$ inputs.
    \item \textbf{PureConvKAN:} three convolutional KAN blocks (16, 32 and 64 channels), each followed by max-pooling and batch normalisation, then global average pooling, Dropout(0.5) and a linear output layer; trained from scratch with $128 \times 128$ inputs.
    \item \textbf{PureKAN:} the image flattened into a vector of pixels and fed directly to a KAN, with no convolution; $64 \times 64$ inputs ($3 \times 64 \times 64 = 12{,}288$ inputs).
\end{itemize}

ResKANet-18 was trained only for entry/exit, and PureConvKAN and PureKAN for both tasks. All were trained for 60 epochs end to end with AdamW and batches of 8 images. They were evaluated on a single 80/20 hold-out split separated by case, which avoids case leakage but is a single split, at lower resolution and with a recipe different from the single protocol. Table~\ref{tab:exploratorios} gives their results.

% Gerado por tcc_dissertation/scripts/gerar_tabelas.py a partir de resultados/. Não edite à mão.
\begin{table}[!htbp]
\centering
\caption{\EP{Exploratory KAN architectures trained from scratch (no ImageNet pre-training) in an earlier phase of this work. They used a single 80/20 hold-out split by case, inputs of 64 or 128 px, batches of 8 and their own recipe, so they are \emph{not} comparable with any number obtained under the single protocol. M = macro, m = micro, W = weighted.}{Arquiteturas KAN exploratórias treinadas do zero (sem pré-treino no ImageNet) numa fase anterior deste trabalho. Usaram um único hold-out 80/20 por caso, entradas de 64 ou 128 px, lotes de 8 e receita própria, portanto \emph{não} são comparáveis com nenhum número obtido no protocolo único. M = macro, m = micro, W = ponderada.}}
\label{tab:exploratorios}
\footnotesize
\setlength{\tabcolsep}{2.5pt}
\resizebox{\textwidth}{!}{%
\begin{tabular}{llcccccccccccccccc}
\toprule
 & & Acc & \multicolumn{3}{c}{\EP{Precision}{Precisão}} & \multicolumn{3}{c}{\EP{Recall}{Revocação}} & \multicolumn{3}{c}{F1} & \multicolumn{3}{c}{AUC} & \multicolumn{3}{c}{\EP{Specificity}{Especificidade}} \\
\EP{Model}{Modelo} & \EP{Input}{Entrada} & & M & m & W & M & m & W & M & m & W & M & m & W & M & m & W \\
\midrule
\multicolumn{18}{l}{\textit{\EP{Entry/exit}{Entrada/saída}}} \\
ResKANet-18 & 128 px & 0.822 & 0.767 & 0.822 & 0.820 & 0.762 & 0.822 & 0.822 & 0.765 & 0.822 & 0.821 & 0.839 & 0.885 & 0.839 & 0.762 & 0.822 & 0.703 \\
PureConvKAN & 128 px & 0.707 & 0.614 & 0.707 & 0.704 & 0.611 & 0.707 & 0.707 & 0.613 & 0.707 & 0.706 & 0.682 & 0.768 & 0.682 & 0.611 & 0.707 & 0.516 \\
PureKAN & 64 px & 0.692 & 0.614 & 0.692 & 0.692 & 0.614 & 0.692 & 0.692 & 0.614 & 0.692 & 0.692 & 0.684 & 0.733 & 0.684 & 0.614 & 0.692 & 0.536 \\
\midrule
\multicolumn{18}{l}{\textit{\EP{MLSD}{MLSD}}} \\
PureConvKAN & 128 px & 0.803 & 0.268 & 0.803 & 0.645 & 0.333 & 0.803 & 0.803 & 0.297 & 0.803 & 0.716 & 0.536 & 0.888 & 0.509 & 0.667 & 0.902 & 0.197 \\
PureKAN & 64 px & 0.712 & 0.355 & 0.712 & 0.689 & 0.324 & 0.712 & 0.712 & 0.328 & 0.712 & 0.693 & 0.463 & 0.841 & 0.496 & 0.679 & 0.856 & 0.326 \\
\bottomrule
\end{tabular}}
\end{table}


The only reading these numbers support is qualitative: without a pre-trained extractor, the KAN architectures reach 0.69--0.82 accuracy in entry/exit, and in the distance task the macro AUC of the two is close to or below chance (0.46--0.54), that is, they learn little about the minority classes. This is consistent with the main finding that the pre-trained backbone, and not the KAN, carries the discrimination.
